% BHU PYQ -- Manually transcribed & error-corrected
% Paper: BCF-123 : Fundamental of Accounting
% Exam: B.Com. (Hons.) FMM Semester II Examination, 2022-23

\documentclass[11pt,a4paper]{article}
\usepackage[a4paper,top=2.5cm,bottom=2.5cm,left=2.8cm,right=2.8cm,headheight=14pt]{geometry}
\usepackage{amsmath,amssymb}
\usepackage[T1]{fontenc}
\usepackage[utf8]{inputenc}
\usepackage{lmodern,microtype,fancyhdr,enumitem,hyperref,lastpage,parskip}

\hypersetup{
    colorlinks=true,
    urlcolor=black,
    linkcolor=black,
    pdftitle={BCF-123: Fundamental of Accounting -- BHU B.Com. (Hons.) FMM Sem II 2022-23}
}

\pagestyle{fancy}
\fancyhf{}
\renewcommand{\headrulewidth}{0.4pt}
\renewcommand{\footrulewidth}{0.4pt}
\fancyhead[L]{\small   \textit{Banaras Hindu University}}
\fancyhead[R]{\small   \textit{BCF-123: Fundamental of Accounting}}
\fancyfoot[C]{\small Page  \thepage\ of \pageref{LastPage}}


\newlist{parts}{enumerate}{2}
\setlist[parts,1]{label=(\alph*),leftmargin=*,itemsep=5pt,topsep=3pt}
\setlist[parts,2]{label=(\roman*),leftmargin=*,itemsep=3pt,topsep=2pt}

\newcommand{\pts}[1]{\hfill   \textit{[#1]}}


\begin{document}


\begin{center}
  {\large\bfseries BANARAS HINDU UNIVERSITY}\\[2pt]
  {
\normalsize B.Com. (Hons.) FMM --- Semester II Examination, 2022-23}\\[6pt]
  \rule{0.8\linewidth}{0.5pt}\\[6pt]
  {\large\bfseries Paper No. BCF-123: Fundamental of Accounting}\\[4pt]
  {
\normalsize Time: 3 Hours\hfill Maximum Marks: 70}
\end{center}

\rule{\linewidth}{0.4pt}

\smallskip



\noindent   \textit{Instructions: Answer all questions. All questions carry equal marks (14 marks each).}

\bigskip




\noindent   \textbf{Question 1.}
What is Accounting? Outline the need for accounting and briefly describe the objects of accounting.\pts{14}

\centerline{   \textbf{OR}}

Pass journal entries for the following transactions in the books of S Ltd. assuming that both parties belong to the same state and CGST @ 6\% and SGST @ 6\% are levied:\pts{14}

\begin{parts}
  \item Purchased goods for Rs. 1,80,000 from Akanksha \& Co.
  \item Sold goods for Rs. 3,50,000 to Nupur Store.
  \item Returned goods to Akanksha \& Co. for Rs. 20,000.
  \item Nupur Store returned goods for Rs. 16,000.
  \item Paid for Printing and Stationery Rs. 10,000.
  \item Goods withdrawn by the proprietor for personal use Rs. 40,000.
  \item Goods destroyed by fire Rs. 30,000.
  \item Payment made of balance of GST.
\end{parts}

\medskip\rule{\linewidth}{0.2pt}


\newpage




\noindent   \textbf{Question 2.}
Trial Balance of a business as at 31st March, 2022 is given below:\pts{14}

\begin{center}
  \small
  
\begin{tabular}{|l|r||l|r|}
    \hline
   \textbf{Debit Balances} &   \textbf{Amount (Rs.)} &   \textbf{Credit Balances} &   \textbf{Amount (Rs.)} \\
    \hline
    Stock on 1st April, 2021 & 25,000 & Sales & 2,27,800 \\
    Furniture & 8,000 & Commission & 500 \\
    Plant and Machinery & 1,500,000 & Returns Outward & 1,000 \\
    Debtors & 30,000 & Creditors & 40,000 \\
    Wages & 12,000 & Capital & 150,000 \\
    Salaries & 20,000 & & \\
    Bad Debts & 1,000 & & \\
    Purchases & 120,000 & & \\
    Electricity Charges & 1,200 & & \\
    Telephone Charges & 2,400 & & \\
    General Expenses & 3,000 & & \\
    Postage Expenses & 1,800 & & \\
    Returns Inward & 900 & & \\
    Insurance Premium & 1,500 & & \\
    Cash in Hand & 2,500 & & \\
    Cash at Bank & 40,000 & & \\
    \hline
   \textbf{Total} &   \textbf{419,300} &   \textbf{Total} &   \textbf{419,300} \\
    \hline
  \end{tabular}
\end{center}
Prepare Trading and Profit and Loss Account for the year ended 31st March, 2022 and Balance Sheet as at that date after taking into account the following adjustments:

\begin{parts}
  \item Closing Stock was valued at Rs. 7,000.
  \item Outstanding liabilities for wages were Rs. 600 and salaries Rs. 1,400.
  \item Depreciation is to be provided @5\% p.a. on all fixed assets.
  \item Included in Plant and Machinery is a machine purchased for Rs. 10,000 on 1st October, 2021.
  \item Insurance premium paid in advance Rs. 200.
\end{parts}

\centerline{   \textbf{OR}}

What are Bills of exchange? Give example. Differentiate between bills of exchange and promissory note.\pts{14}

\medskip\rule{\linewidth}{0.2pt}


\newpage




\noindent   \textbf{Question 3.}

\begin{parts}
  \item A and B share profits in the ratio of A: $5/8$ and B: $3/8$. C is admitted as a partner. He brings in Rs. 70,000 as his capital and Rs. 48,000 as goodwill. The new profit-sharing ratio among A, B and C is agreed to be 7 : 5 : 4 respectively. Pass Journal entries.\pts{7}
  \item X and Y were partners sharing profits in the ratio of 5 : 3 respectively. On 1st April, 2022 they admitted Z as a new partner, all the partners agreeing to share future profits equally. On the date of admission of the new partner, there was a goodwill account in the old firm's ledger showing a balance of Rs. 18,000. The current value of firm's goodwill was placed at Rs. 36,000. Z paid Rs. 50,000 by way of his capital. He also paid an appropriate amount for his share of goodwill. X and Y have written-off the goodwill account before Z's admission. Pass Journal entries.\pts{7}
\end{parts}

\centerline{   \textbf{OR}}

X and Y were partners sharing profits in the ratio of 5 : 3 respectively. The Balance Sheet of the partnership firm as on 31st March, 2022 was as follows:\pts{14}

\begin{center}
  \small
  
\begin{tabular}{|l|r||l|r|}
    \hline
   \textbf{Liabilities} &   \textbf{Amount (Rs.)} &   \textbf{Assets} &   \textbf{Amount (Rs.)} \\
    \hline
    Capital Accounts: & & Goodwill & 190,000 \\
    \quad X & 205,000 & Building & 85,000 \\
    \quad Y & 165,000 & Machinery & 54,740 \\
    Profit \& Loss Account & 56,000 & Sundry Debtors & 30,000 \\
    Sundry Creditors & 27,400 & Stock & 72,630 \\
    & & Cash at Bank & 21,030 \\
    \hline
   \textbf{Total} &   \textbf{453,400} &   \textbf{Total} &   \textbf{453,400} \\
    \hline
  \end{tabular}
\end{center}
On the above date, Z was admitted on the following terms:

\begin{parts}
  \item Z would get 1/5th share in the profits.
  \item Z would pay Rs. 1,20,000 as capital and Rs. 16,000 for his share of goodwill.
  \item Machinery would be depreciated by 10\% and building would be appreciated by 30\%. A provision for bad debts @ 5\% on debtors would be created. An unrecorded liability amounting to Rs. 3,000 for repairs to building would be recorded in the books of account.
  \item Immediately after Z's admission, goodwill account would be written off. Thereafter, the capital accounts of the old partners would be adjusted through the necessary current accounts in such a manner that the capital accounts of all the partners would be in their profit sharing ratio.
\end{parts}
Prepare Revaluation account, Capital accounts and the Initial Balance Sheet of the new firm.

\medskip\rule{\linewidth}{0.2pt}


\newpage




\noindent   \textbf{Question 4.}
At the time of admission of a partner, why must assets and liabilities be revalued? Explain and advise how it is treated in the books of account.\pts{14}

\centerline{   \textbf{OR}}

D, K and N are partners sharing profit and loss in the ratio of 2 : 2 : 1. N retires on 31-3-2022. The Balance sheet of the firm as on 31-3-2022 was as under:\pts{14}

\begin{center}
  \small
  
\begin{tabular}{|l|r||l|r|}
    \hline
   \textbf{Liabilities} &   \textbf{Amount (Rs.)} &   \textbf{Assets} &   \textbf{Amount (Rs.)} \\
    \hline
    Capital Accounts: & & Goodwill & 10,000 \\
    \quad D & 30,000 & Machinery & 20,000 \\
    \quad K & 20,000 & Investment & 10,000 \\
    \quad N & 10,000 & Sundry Debtors & 30,000 \\
    General Reserve & 5,000 & Stock & 10,000 \\
    Investment Fluctuation Fund & 2,500 & Cash at Bank & 5,000 \\
    Bad Debts Reserve & 2,000 & & \\
    Creditors & 15,500 & & \\
    \hline
   \textbf{Total} &   \textbf{85,000} &   \textbf{Total} &   \textbf{85,000} \\
    \hline
  \end{tabular}
\end{center}
Following adjustments are agreed at the time of retirement:

\begin{parts}
  \item Value of machinery will be Rs. 25,000 and Value of stock is Rs. 5,000.
  \item Value of investments is Rs. 8,000, which is taken by N at this price.
  \item An amount of Rs. 5,000 included in creditors is no longer payable.
  \item The provision for workmen compensation to be credited at Rs. 2,000.
  \item The provision for doubtful debts is to be kept at 10\% on debtors.
  \item Goodwill of the firm is valued at Rs. 40,000.
\end{parts}
Prepare necessary accounts and the balance sheet of the firm after N's retirement.

\medskip\rule{\linewidth}{0.2pt}




\noindent   \textbf{Question 5.}
How the Joint Life Policy is recorded at the time of retirement of a partner?\pts{14}

\centerline{   \textbf{OR}}

How will you calculate the amount payable to the executors of a deceased partner?\pts{14}

\vfill
\rule{\linewidth}{0.4pt}

\begin{center}\small
\href{https://bhu-exam.vercel.app/}{Click here to visit our website}\\[4pt]
\href{https://me-aryan.vercel.app/}{Click here to visit our contributor}\\[4pt]
\href{https://quashan-me.vercel.app/}{Click here to visit our contributor}
\end{center}

\end{document}
